\documentclass[11pt]{article}
\usepackage{mas1004}

\reporttitle{Assignment 2: Long report}
\studentname{Kim Yunseo}
\studentid{20251129}

\begin{document}
\makeheader

\section*{Track A}

\begin{center}\footnotesize
\begin{tabular}{lp{0.55\linewidth}r}
\toprule
tag & what changed from the run before & validation \\
\midrule
run 0 & Starter baseline recipe (Crop scale 0.35--1.0, ColorJitter 0.2) & 0.4956 \\
run 1 & Crop scale 0.5--1.0, ColorJitter 0.3, added GaussianBlur & 0.4956 \\
run 2 & Removed GaussianBlur, relaxed ColorJitter to 0.15 (scale 0.45--1.0) & 0.4956 \\
\bottomrule
\end{tabular}
\end{center}

\subsection*{Experimental Design and Observations}

In Track A, the goal was to classify street-view images from five cities (Seoul, Busan, Tokyo, Osaka, Taipei) using a fixed ResNet18 model within a strict budget of 249,984 training images. We followed the principle of changing one main variable at a time to systematically isolate the effect of image augmentations.

\begin{itemize}
    \item \textbf{Run 0 (Baseline)}: The starter setup achieved an overall validation accuracy of 49.56\%. However, there was a severe performance imbalance across cities: Tokyo scored 62.82\% and Seoul scored 57.00\%, whereas Osaka (34.02\%) and Busan (38.56\%) performed very poorly. The model appeared overfitted to clean images with clear structural landmarks.
    \item \textbf{Run 1 (Exp 2 - Strong Augmentation \& Blur)}: To improve Osaka's low accuracy, we increased the crop scale lower bound to 0.5 to prevent losing entire landmarks, strengthened \texttt{ColorJitter} (0.3), and introduced \texttt{GaussianBlur}. Osaka accuracy jumped significantly from 34.02\% to 40.27\% (+6.25\%p). However, Tokyo accuracy dropped from 62.82\% to 59.44\%, leaving the overall score unchanged at 49.56\%.
    \item \textbf{Run 2 (Exp 3 - Balanced Augmentation)}: Hypothesizing that Gaussian blur destroyed crisp edge features crucial for Tokyo and Seoul (such as road signs and crisp building text), we removed the blur and relaxed color jitter. As expected, Tokyo (61.99\%) and Seoul (57.78\%) recovered, but Osaka dropped back to 36.94\%, yielding 49.56\% once again.
\end{itemize}

\subsection*{Final Recipe and Validation vs. Leaderboard Score}

Our final submitted recipe (\texttt{track\_a/recipe.py}) utilizes the balanced augmentation strategy from Run 2 along with the full selection dataset (\texttt{track\_a/selection.csv}).

\begin{center}
\begin{tabular}{lrr}
\toprule
Metric & Validation Score & Leaderboard Score \\
\midrule
Overall Accuracy & 0.4956 & 0.4956 \\
\bottomrule
\end{tabular}
\end{center}

\section*{Decisions Made from Photos}

\begin{center}
\small
\begin{tabular}{p{0.45\linewidth}p{0.45\linewidth}}
\toprule
\textbf{Tokyo / Seoul Features} & \textbf{Osaka / Busan Features} \\
\midrule
Clear visual landmarks, bright daylight lighting, sharp text/signboard edges. & Higher proportion of night/dark lighting, wide angles, and environmental noise. \\
\bottomrule
\end{tabular}
\end{center}

When inspecting the training photos, we observed that images from Tokyo and Seoul generally contained sharp lighting, clear signboard text, and distinct architectural features. In contrast, many images from Osaka and Busan were darker, captured from wider/blurry angles, or contained high environmental noise.

Before taking action, we established a rule: \textit{"Do not apply heavy spatial or Gaussian blur augmentations if the city's key features rely on fine text and crisp edges, unless dataset-level cleaning is performed."}

Because modifying image transforms alone (\texttt{recipe.py}) consistently hit a wall at 49.56\% due to trade-offs between city characteristics, any further improvement requires dataset-level curation (\texttt{selection.csv}) to balance class counts and filter out non-street photos (such as indoor scenes or documents).

\section*{What Made the Difference}

\begin{itemize}
    \item \textbf{Most Impactful Change}: Adding \texttt{GaussianBlur} and stronger \texttt{ColorJitter} in Run 1 produced the single largest per-city gain (+6.25\%p in Osaka), proving that regularization helps low-quality, noisy images.
    \item \textbf{Bottleneck}: Modifying data augmentations alone without cleaning or rebalancing the underlying dataset yielded zero net change in overall accuracy across all three runs (stuck at 49.5
